%% ECON 282E -- Final project guidelines ("Assemble"), written 2026-09-26.
%% Public. The final project assembles components each student built in HW1 and HW2.
%% Adapted in part from AI-ECON-2026/source/Docs/Capstone_Projects_2026.md (shared rules,
%% deliverables, rubric), without that course's tracks or any unpublished research.
%% Build:  python3 tools/build_docs.py homeworks-exams/FinalProject/guidelines.tex

\documentclass[11pt]{article}

\usepackage[margin=1in]{geometry}
\usepackage{amsmath,amssymb,booktabs,enumitem,xcolor,tabularx,array}
\usepackage[T1]{fontenc}
\usepackage{microtype}
\usepackage[colorlinks=true,linkcolor=blue!50!black,urlcolor=blue!50!black]{hyperref}

\newcolumntype{L}{>{\raggedright\arraybackslash}X}

\title{\vspace{-2em}\textbf{ECON 282E --- Final Project Guidelines}\\[2mm]
       \large Assemble: From Components to a Research Answer}
\author{Foundations of Macroeconomics \\ University of California, Riverside}
\date{Teams: Thursday, November 5 \quad\textbullet\quad Talks: Thursday, December 3
      \quad\textbullet\quad Paper: Thursday, December 10, 2026}

\begin{document}
\maketitle
\thispagestyle{empty}

\noindent\rule{\linewidth}{0.4pt}\par
\small
\textbf{The idea.} In this course you learn to be the \textbf{project manager} of a quantitative
research project. You formalize the question, choose the model, define the equilibrium, pick the
method, find the data, check quality, and communicate the answer. AI assistants write much of the
code; the decisions and the validation are yours. The two homework assignments each build one
\emph{component} of that skill set. The final project \textbf{assembles} them, in a team, to answer
a question the team poses itself.

\textbf{Weight.} 40\% of the course grade: proposal 5\%, talk and discussant role 15\%, paper and
replication package 20\%.
\normalsize\par
\noindent\rule{\linewidth}{0.4pt}\par

%=============================================================================
\section{The ladder: HW1 \texorpdfstring{$\to$}{->} HW2 \texorpdfstring{$\to$}{->} project}

Each rung removes scaffolding and adds ownership.

\medskip
{\small
\renewcommand{\arraystretch}{1.25}
\noindent\begin{tabularx}{\linewidth}{@{}p{2.6cm}LLL@{}}
\toprule
 & \textbf{HW1: the representative-agent engine} & \textbf{HW2: heterogeneity and learning} & \textbf{Final project: assembly} \\
\midrule
Sessions & 1--5 & 6--9 & all; Session 10's agentic workflow \\
Dates & out Oct 8 and 15; due \textbf{Oct 29} & out Oct 29 and Nov 12; due \textbf{Mon Nov 23} & teams Nov 5; proposal Nov 16; talk Dec 3; paper Dec 10 \\
Who sets the question & the instructor (common spine), plus \textbf{your own pitch} (B5) & a shared core, plus \textbf{a route question you choose} & \textbf{your team}, evolved from the pitches \\
What you learn to manage & charter, unit test, data provenance, method choice, reproducibility & scoping a route, benchmarking an ML solver against a classical one, auditing someone else's work & integration, division of labour, internal quality control, discussing another team's work \\
Quality control & AI red-team item; validation ledger & AI red-team item; \textbf{peer replication audit} of a classmate's HW1 & internal audit; discussant report \\
What you carry forward & a calibrated RA model, a data pipeline, a method library, a pitch & an HA solver and one learning or measurement module & --- \\
\bottomrule
\end{tabularx}
}

%=============================================================================
\section{The integration rule}

The team's question must be answered by \textbf{combining at least one HW1 component and at least
one HW2 component built by its members}, plus the piece the team adds: a new model feature, new
data, or a new method. Examples of components:
\begin{itemize}[leftmargin=*,itemsep=1pt]
  \item \textbf{From HW1:} the FRED data and calibration pipeline (A3); a validated global,
        perturbation or collocation solver (B2); an estimation routine (route R2); a capital-tax
        extension (route R3); a second-order solution (route R1).
  \item \textbf{From HW2:} the Aiyagari solver and stationary distribution (C1); the Euler-residual
        network (C2); an RL agent (C3); a text-based index (C4); a Krusell--Smith solver (route H1);
        a transition-path solver (route H2); a deep or RL solver for the HA household (route H3); a
        text- or exposure-disciplined income process (route H4).
\end{itemize}

\textbf{Validate before you extend.} Each reused component must first reproduce its own homework
benchmark inside the team's repository --- the same numbers, from the team's code, on the team's
machine --- before it is extended. This is the course's version of ``replicate first, explore
second'': the trusted benchmark comes before the new result. A project that assembles its
components correctly and validates them rigorously, with a modest extension, earns a strong grade.

%=============================================================================
\section{Required elements}

Every project, whatever its emphasis, must have:
\begin{enumerate}[leftmargin=*,itemsep=2pt]
  \item \textbf{A question the team formalized}, stated in one sentence, and labelled positive or
        normative. Say what number or ranking answers it.
  \item \textbf{A model} from the Session 2 and Session 6 family (Aiyagari, Krusell--Smith, OLG,
        optimal policy), or a justified alternative. Say why a simpler model will not do, and define
        the equilibrium.
  \item \textbf{One classical and one AI/ML method, cross-validated on at least one object} (a policy
        function, a price, a distribution, an index). Report where they disagree and why.
  \item \textbf{Data with provenance}, structured (FRED, BLS, PSID, \dots) or text, acquired by
        script, with series IDs or URLs and access dates. Public data only.
  \item \textbf{An agentic workflow} (Session 10). The repository carries a \texttt{CLAUDE.md} (or the
        equivalent file for your agent) that describes the project to the next agent and the next
        human, and an \texttt{AI\_LOG.md} that records the prompts that mattered, the decisions you
        made, where the AI was wrong and how you caught it, and an honest split of human and AI
        contribution.
  \item \textbf{An internal audit}: each member re-runs a component built by another member and
        records the result with \texttt{templates/audit\_report.md}.
  \item \textbf{One-command replication}: \texttt{python run\_all.py} (or equivalent) reproduces every
        number and figure from raw inputs. Fixed seeds; pinned environment; a \texttt{README} with
        run instructions. \textbf{No API keys in the repository} (use \texttt{.env} and
        \texttt{.gitignore}).
\end{enumerate}

\textbf{Never submit code you have not read.} Every member must be able to explain every result in
the paper.

%=============================================================================
\section{Emphases}

The syllabus lists four tracks. They are not separate projects but \textbf{emphases}: they say
where the team's added piece lies. Each connects to homework routes whose components it naturally
reuses.

\medskip
{\small
\renewcommand{\arraystretch}{1.2}
\noindent\begin{tabularx}{\linewidth}{@{}p{4.4cm}Lp{3.6cm}@{}}
\toprule
\textbf{Emphasis} & \textbf{The added piece} & \textbf{Natural components} \\
\midrule
\textbf{A.} A quantitative model and a policy experiment & a new friction, policy or counterfactual in an HA or OLG model, with welfare by group & HW1 R3; HW2 C1, H1, H2 \\
\textbf{B.} A hard model solved with deep learning or RL & a learning solver where grids strain, validated against a classical benchmark & HW1 R1; HW2 C2, C3, H3 \\
\textbf{C.} The economics of AI & AI as a shock to tasks, wages or capital, in a model you can solve & HW1 R2 or R3; HW2 C1, H4 \\
\textbf{D.} Text or LLM measurement feeding a model & a validated text-based measure that disciplines a parameter or a shock & HW1 R2; HW2 C4, H4 \\
\bottomrule
\end{tabularx}
}

\medskip
Teams may mix emphases. A strong project usually crosses one line of this table: for example, a
text-measured shock (D) fed into an Aiyagari policy experiment (A).

%=============================================================================
\section{Milestones}

\begin{center}
\small
\renewcommand{\arraystretch}{1.2}
\begin{tabularx}{\linewidth}{@{}p{3.3cm}L@{}}
\toprule
\textbf{Date} & \textbf{Milestone} \\
\midrule
Thu Oct 29 & \textbf{HW1 due, including your one-page pitch (B5).} Post the pitch, under your name or
anonymously, in the forum category \emph{Pitches and teams}, and read the others. \\
Thu Nov 5 & \textbf{Teams (2--3) and topic}, posted in \emph{Pitches and teams}: members, the one-sentence
question, and which members' pitches it grew from. \\
Mon Nov 16 & \textbf{Two-page proposal} (5\%; \texttt{templates/proposal.md}): the question; the HW1 and HW2
components you will reuse and their benchmarks; the gap your team fills; the division of
labour; a risk register. Written feedback by Session 9 (Nov 19). \\
Mon Nov 23 & \textbf{HW2 due.} Your HW2 components should now be ready to reuse. \\
Thu Dec 3, Block B & \textbf{Talks and discussants} (15\%). Slots: with at most five teams, a 10-minute talk,
3 minutes from the discussant and 2 of Q\&A; with six teams, 8 + 2 + 2. Each team discusses
another team's project (\texttt{templates/discussant\_form.md}); the pairing and the other team's
slides are sent by Monday, November 30. \\
Thu Dec 10 & \textbf{Paper} of at most 10 pages, \textbf{replication package}, \texttt{AI\_LOG.md}, internal
audit reports and the discussant form (20\%). \\
\bottomrule
\end{tabularx}
\end{center}

%=============================================================================
\section{Deliverables}

\begin{enumerate}[leftmargin=*,itemsep=2pt]
  \item \textbf{Paper} (at most 10 pages, references included): question, model and equilibrium,
        components reused and how each was validated, calibration and data provenance, method and
        its accuracy, results, limitations. Written like a short working paper. Every number in it
        must appear in the output of the replication command.
  \item \textbf{Repository}: \texttt{CLAUDE.md}, \texttt{README}, pinned environment, fixed seeds,
        data-acquisition scripts (no manual downloads), a short data dictionary, and one-command
        reproduction.
  \item \textbf{\texttt{AI\_LOG.md}}, with a validation ledger linking each headline claim to the check
        that supports it, as in HW1.
  \item \textbf{Internal audit reports}, one per member.
  \item \textbf{Talk slides} and the \textbf{discussant form} for the team you discussed.
\end{enumerate}

%=============================================================================
\section{Rubric}

\begin{center}
\small
\renewcommand{\arraystretch}{1.2}
\begin{tabularx}{\linewidth}{@{}p{4.6cm}rL@{}}
\toprule
\textbf{Dimension} & \textbf{Points} & \textbf{What we look for} \\
\midrule
Framing and question & 15 & A precise, well-motivated question; the right model and equilibrium concept; why a simpler model will not do \\
Integration and correctness & 20 & Reused components reproduce their benchmarks; they are combined correctly; the added piece is right \\
Validation & 20 & Hand checks, benchmarks, a classical-vs-ML cross-validation, sanity checks; honest about what failed \\
Reproducibility and provenance & 15 & One-command reproduction on a clean machine; scripted data with access dates; pinned environment \\
AI-collaboration quality & 10 & An architect, not a consumer: an informative \texttt{AI\_LOG.md}, a useful \texttt{CLAUDE.md}, honest attribution \\
Communication & 10 & A clear paper and talk; figures that tell the story \\
Discussant and audit & 10 & A specific, fair, well-evidenced discussant report; internal audits that actually re-ran the code \\
\bottomrule
\end{tabularx}
\end{center}

\textbf{How the rubric maps to the course grade.} The proposal (5\%) is graded on the framing
dimension and on the quality of the component plan and risk register. The talk and discussant role
(15\%) is graded on communication and on the discussant dimension. The paper and replication
package (20\%) are graded on the remaining dimensions. All members of a team receive the same
project grade unless the audit reports or the \texttt{AI\_LOG.md} show a clearly unequal
contribution.

%=============================================================================
\section{Coverage: where each session's tools are used}

The two homework assignments and the project together exercise every key model and technique of
the course. \textbf{Core} items are required of everyone; \textbf{route} items deepen one area by
choice; the project \textbf{integrates} them.

\medskip
{\small
\renewcommand{\arraystretch}{1.2}
\noindent\begin{tabularx}{\linewidth}{@{}p{0.95cm}Lp{2.3cm}p{2.3cm}p{2.6cm}@{}}
\toprule
\textbf{Sess.} & \textbf{Key models and techniques} & \textbf{HW1} & \textbf{HW2} & \textbf{Project} \\
\midrule
1 & growth model; planner and equilibrium; dynamic programming; why RL & A1, A2 (core) & C3 & model choice \\
2 & Bewley--Huggett--Aiyagari; Krusell--Smith; OLG; Ramsey and time consistency & A1(d) (core); R3 & C1 (core); H1, H2 & emphasis A, C \\
3 & VFI; Euler equations and time iteration; Python, PyTorch; data workflow & A2, A3 (core) & C2 & replication package \\
4 & root-finding, optimization, interpolation, quadrature; Tauchen and Rouwenhorst & B1, B2 (core); R2 & C1 & calibration, estimation \\
5 & perturbation, Blanchard--Kahn, pruning; Chebyshev collocation, Smolyak, finite elements & B2 (core); R1 & C2 benchmark & emphasis B \\
6 & Aiyagari as a nested fixed point; Krusell--Smith; OLG transitions; Negishi & --- & C1 (core); H1, H2 & emphasis A \\
7 & neural networks; training; Euler-residual and deep VFI solvers; validation & --- & C2 (core); H3 & classical-vs-ML check \\
8 & Q-learning, policy gradients, actor--critic; HA at the frontier; taxation without commitment & --- & C3 (core); H1, H2, H3 & emphasis B \\
9 & text as data; embeddings; LLMs as instruments; retrieval & A3(c) (AI audit) & C4 (core); H4 & emphasis C, D \\
10 & agentic research workflows; verification; governance & AI red-team & AI red-team & \texttt{CLAUDE.md}, \texttt{AI\_LOG.md}, audit \\
\bottomrule
\end{tabularx}
}

\vfill
\noindent\rule{\linewidth}{0.4pt}\par
\small
\textbf{Ground rules.} AI assistance is encouraged as a partner, not a ghostwriter. Public data
only; no personal, IRB- or HIPAA-protected data. Cite every source and separate your contribution
from prior work. Templates: \texttt{homeworks-exams/FinalProject/templates/}.
\normalsize

\end{document}
